\section{Simulating an arbitrary discrete distribution using unbiased coin tosses}

We can generalize Simulation Method 2 described above
to get a method for simulating an arbitrary discrete
random variable taking values
$\alpha_1,\ldots,\alpha_k$ with respective
probabilities $p_1,\ldots,p_k$, using a sequence of
independent unbiased coin tosses
$X_1,X_2,\ldots \sim \operatorname{Ber}(1/2)$. Let
$U=\sum_{n=1}^\infty X_n/2^n$ as before, and for
$m\ge 1$ let $U_m = \sum_{k=1}^m X_k/2^k$ be the
partial sums of the binary expansion of $U$.

$\textbf{Simulation Method 3.}$ Denote
$c_0 = 0, c_j = p_1+\ldots+p_j$. Note that the
intervals $(c_0,c_1), (c_1,c_2),\ldots, (c_{k-1},c_k)$
form a partition of the interval $(0,1)$ into
subintervals of lengths $p_1,\ldots,p_k$. To produce
a sample from the discrete distribution with atoms of
size $p_j$ at $\alpha_j$ for $j=1,\ldots,k$:

\begin{enumerate}[label=\arabic*.]
\item Sample $X_1,X_2,\ldots$ one at a time until the first
time $m$ for which there is a $j$ such that the
condition

\[
c_{j-1} < U_m < U_m+\frac{1}{2^m} < c_j
\]

\noindent holds. Note that this condition can be checked by
comparing the binary string $(X_1,\ldots,X_m)$
against the first $m$ digits in the binary expansions
of the numbers $c_0,\ldots,c_k$.

\item Output $\alpha_j$.
\end{enumerate}

\begin{xexercise}
Show that the condition
$c_{j-1} < U_m < U_m+\frac{1}{2^m} < c_j$ is
equivalent to $c_{j-1} < U < c_j$, and that therefore
for each $1\le j\le k$ the output $\alpha_j$ is
obtained with probability $p_j$.
\end{xexercise}

\noindent The average number of samples from $X_1,X_2,\ldots$
required for the simulation looks like a highly
nontrivial function of the probabilities
$p_1,\ldots,p_k$. Amazingly, to a good approximation
it is equal to the $\textbf{entropy}$ of the
distribution $p_1,\ldots,p_k$, a well-known function
from information theory that is known to measure the
information content of a random variable. The precise
result to that effect is as follows.

\begin{thm}
The $\textbf{entropy}$ of the discrete probability
distribution $p_1,\ldots,p_k$ is defined by
$H(p_1,\ldots,p_j) = -\sum_{j=1}^k p_j \log_2 p_j$.
The number of samples $N$ required in the simulation
satisfies

\[
H(p_1,\ldots,p_k) \le \mathbf{E} N \le H(p_1,\ldots,p_k) + 4.
\]

\end{thm}

\noindent For the proof, see my paper \emph{Sharp entropy bounds for
discrete statistical simulation (Stat. Prob. Lett. 42
(1999), 219--227)}, or section 5.11 of the book
\emph{Elements of Information Theory, 2nd Ed.} by T. M.
Cover and J. A. Thomas, where a similar result is
proved for a different simulation method for which
$\mathbf{E} N$ can be bounded from above by the
slightly better bound $H(p_1,\ldots,p_k)+2$.
